\section{How to Evolve}
\label{sec:how}

\begin{table*}[!t]
\centering
\caption{Overview of Reward-based, Imitation/Demonstration, and Population-based Learning Methods for Self-Evolving Agents. This table categorizes key approaches based on the following criteria: (1) Feedback Type: the type of feedback used, including language-based rationales and numerical rewards. (2) Feedback Source: the origin of the feedback, either internal (model-generated) or external (provided externally). (3) Learning Method: the learning paradigm applied, such as in-context learning (ICL), supervised fine-tuning (SFT), reinforcement learning (RL), and evolutionary algorithms; (4) Updated Components: which parts of the model are updated, either full parameters or a subset of the model. (5) Update Timing: the stage during the agent's evolution when updates are applied, such as pre-training, pre-test, or test-time.}
\resizebox{\textwidth}{!}{
\begin{tabular}{ccccccc}
\toprule
\textbf{Method} & \textbf{Feedback Type} & \textbf{Feedback Source} & \textbf{Learning Method} & \textbf{Updated Components} & \textbf{Update Timing}  \\
\midrule
\multicolumn{6}{c}{\textit{Reward-based Evolution Methods}} \\\midrule
Reflexion\citep{shinn2023reflexion} & language & internal & ICL & context & test-time \\
AdaPlanner\citep{sun2023adaplanner} & language & external + internal & ICL & context & test-time \\
AgentS2\citep{agashe2025agents2} & language & external & ICL & context & test-time \\
SELF\citep{lu2023self} & language & external + internal & SFT & full params & pre-test time + test-time \\
SELF-REFINE\citep{madaan2023self_refine} & language & internal & ICL & context & test-time \\
SCoRe\citep{kumar2024training} & numerical & external  & RL & full params & pre-test time \\
PAG\citep{jiang2025pag} & numerical & external  & RL & full params & pre-test time \\
TextGrad\citep{yellamraju2024textgrad} & language & external & ICL & context & pre-test time / test-time \\
SRSI\citep{simonds2025selfrewardingselfimproving} & language & internal & RL & full params & pre-test time \\
Self-Train LM\citep{shafayat2025largereasoningmodelsselftrain} & numerical & internal & RL & full params & pre-test time \\
MM-UPT\citep{wei2025unsupervisedposttrainingmultimodalllm} & numerical & internal & RL & full params & pre-test time \\
CoVo\citep{zhang2025consistentpathsleadtruth} & numerical & internal & RL & full params & pre-test time \\
SWE-agent\citep{du2025swedevevaluatingtrainingautonomous} & language & external & ICL & context & test-time \\
SICA\citep{robeyns2025self} & numerical & external & ICL & codebase(tools, workflows, prompts) & test-time \\
Feedback Friction\citep{jiang2025feedbackfriction} & language & external & ICL & context & test-time \\
USEagent\citep{applis2025useagent} & language & external & ICL & context & test-time \\
DYSTIL\citep{wang2025dystildynamicstrategyinduction} & language + numerical & external + internal & SFT+RL & full params & pre-test time + test-time \\
OTC-PO\citep{wang2025otcpo} & numerical & external & RL & full params & pre-test time \\
AUTORULE\citep{wang2025autorule} & language + numerical & external + internal & RL & full params & pre-test time \\
EGSR\citep{zhang2025learning} & numerical & external  & RL & full params & pre-test time \\
LADDER\citep{simonds2025ladder} & numerical & external  & RL & full params & pre-test time \\
RAGEN\citep{wang2025ragen} & numerical & external & RL & full params & test-time \\
SPIRAL\citep{liu2025spiral} & numerical & internal & RL & full params & pre-test time \\
ICRL Prompting\citep{song2025reward} & numerical & external + internal & RL & full params & test-time \\
MATH-SHEPHERD\citep{wang2023math} & numerical & external & RL & full params & pre-test time \\
AgentPRM\citep{choudhury2025process} & numerical & external & SFT+RL & full params & pre-test time \\
Agent Q\citep{putta2024agentq} & numerical & external & RL & full params & pre-test time \\
GiGPO\citep{feng2025group} & numerical & external & RL & full params & pre-test time \\
SPA-RL\citep{wang2025spa} & numerical & external & RL & full params & pre-test time \\
Self-Instruct\citep{wang2022self} & language & internal & SFT & full params & pre-test time \\
WizardLM\citep{xu2024wizardlm} & language & internal & SFT & full params & pre-test time \\
OS-Genesis\citep{sun2024genesis} & numerical & external & SFT & full params & pre-test time \\
UI-Genie\citep{xiao2025ui} & numerical & external & SFT & partial params & pre-test time \\
GUI-R1\citep{luo2025gui} & numerical & external & SFT+RL & full params & pre-test time \\
InfiGUI-R1\citep{liu2025infigui} & numerical & external & SFT+RL & full params & pre-test time \\
Voyager\citep{wang2023voyager} & language & external & ICL & context & test-time \\
SwiftSage\citep{lin2023swiftsage} & language & external & ICL & context & test-time \\
AutoWebGLM\citep{lai2024autowebglm} & language & external & SFT+RL & full params & pre-test time \\
DigiRL\citep{bai2024digirl} & language & external & RL & partial params & pre-test time \\
WebRL\citep{qi2024webrl} & language & external & SFT+RL & full params & pre-test time \\
Let's Verify Step-by-Step\citep{lightman2023let} & language & external & SFT & full params & pre-test time \\
AlphaMath\citep{chen2024alphamath} & numerical & external & SFT & full params & pre-test time \\
rStar-Math\citep{guan2025rstar} & numerical & external & SFT & full params & pre-test time \\
DistRL\citep{wang2024distrl} & language & external & RL & full params & pre-test time + test-time \\
MobileGUI-RL\citep{shi2025mobileguirladvancingmobilegui} & language & external & RL & full params & pre-test time \\ \midrule
\multicolumn{6}{c}{\textit{Imitation and Demonstration Learning Methods}} \\\midrule
STaR\citep{zelikman2022star} & language + numerical & internal & SFT & full params & pre-test time  \\
V-STaR\citep{hosseini2024v} & numerical & external + internal & SFT + RL & partial params & pre-test time \\
AdaSTaR\citep{koh2025adastar} & numerical & internal & SFT & full params  & pre-test time \\
STIC\citep{deng2024enhancing} & language & internal & RL + SFT & partial params & pre-test time \\
GENIXER\citep{zhao2024genixer}& language & external & SFT & full params & pre-training  \\
SiriuS\citep{zhao2025sirius} & language + numerical & internal & SFT & full params & pre-test time \\
SOFT\citep{tang2025bridging} & language & internal & SFT & not specified & pre-test time \\
RISE\citep{qu2024recursive} & language + numerical & internal + external & SFT & full params & pre-test time  \\
IoE\citep{li2024confidence} & numerical & internal & / & / & test-time \\\midrule
\multicolumn{6}{c}{\textit{Population-based and Evolutionary Methods}} \\\midrule
DGM\citep{zhang2025darwin} & numerical & external & ICL & codebase (tools, workflows, prompts) & test-time \\
EvoMAC\citep{hu2024self} & language & external & ICL & team composition, workflow, prompts & test-time \\
SPIN\citep{chen2024self} & language & internal & RL & full params & pre-test time \\
GENOME\citep{zhang2025nature} & numerical & external & Evolution Alg. & partial params & pre-test time \\
SPC\citep{chen2025spc} & numerical & internal & SFT+RL & critic params & pre-test time + test-time \\
Puppeteer\citep{dang2025multi} & numerical & external & RL & planner policy & pre-test time / between tasks \\
MedAgentSim\citep{almansoori2025self} & language & external & ICL & context (knowledge base) & test-time \\
STL\citep{mendes2025language} & language + numerical & internal & SFT & value model & pre-test time \\
MDTeamGPT\citep{chen2025self} & language & external & ICL & context (knowledge base) & test-time \\
\bottomrule
\end{tabular}
}
\label{tab:imitation_learning}
\end{table*}






The pursuit of self-evolution lies at the heart of building advanced, autonomous, and increasingly general artificial intelligence. For large language models (LLMs) and their agentic extensions, the question of how to continually, autonomously, and efficiently evolve their capabilities has become a central challenge. Therefore, the third key aspect of a self-evolving agent is to instantiate an effective \textit{evolving strategy} $f$, i.e., how to transform an agent system $\Pi=(\Gamma,\{\psi_i\}, \{C_i\}, \{\mathcal{W}_i\})$ to its new state $\Pi'=(\Gamma',\{\psi'_i\}, \{C'_i\}, \{\mathcal{W}'_i\})$. Unlike traditional approaches that rely on static datasets or one-time supervised fine-tuning, self-evolution emphasizes an ongoing process where models learn from real-world interactions, actively seek feedback, self-reflect, generate or curate new data, and adapt their strategies in response to dynamic environments. This continuous evolution is not merely a matter of scaling up data or computation; it requires the agent to acquire a spectrum of meta-capabilities, including self-correction, autonomous data generation, knowledge transfer, and multi-agent collaboration. As a result, the landscape of self-evolution has become increasingly rich and multi-faceted, with each methodological branch exploring different axes of feedback, learning paradigms, data sources, and evolutionary scales.


Over time, research on self-evolving agents has progressed through three major paradigms—reward-based, imitation-based, and population-based evolution—each emerging to address the limitations of the previous one. 
Reward-based methods first closed the feedback loop through explicit signals but suffered from brittleness and high cost. 
Imitation-based learning stabilized evolution by leveraging high-quality demonstrations, though sometimes at the expense of exploration. 
Population-based evolution then extended adaptation to collective scales, emphasizing diversity and emergent coordination. 
Together, these paradigms outline a coherent trajectory from individual self-improvement toward collective intelligence.



This chapter aims to systematically map and analyze the major families of self-evolution methods, providing a unified framework for understanding their principles, mechanisms, and interactions. We begin with \textbf{reward-based evolution}, which centers on the design of reward signals—ranging from natural language feedback and internal confidence metrics to external or implicit signals—to guide iterative self-improvement. 

Next, we examine \textbf{imitation and demonstration learning}, where agents learn by mimicking complete, high-quality behavioral exemplars (i.e., demonstrations). While traditionally sourced from human experts, in the context of self-evolving agents, these exemplars are often generated by the agent itself or by other agents.
This paradigm is particularly powerful when demonstrations are abundant or can be autonomously synthesized, and it has driven significant progress in both reasoning and multimodal domains. 

Finally, we introduce \textbf{population-based and evolutionary methods}, which draw inspiration from biological evolution and collective intelligence. These approaches maintain populations of agent variants or collaborating agents, leveraging mechanisms such as selection, mutation, crossover, and competition to explore the solution space in parallel, foster diversity, and enable the emergence of novel strategies or architectural innovations.



\subsection{Reward-based Self-Evolution}

The capacity for self-improvement is a cornerstone of advanced intelligence. In the context of Large Language Models (LLMs), this manifests as a dynamic process of reward-driven evolution, where models iteratively learn from their own outputs and interactions to refine their capabilities. The design of the reward signal, which serves as the guiding feedback, is crucial; it determines the nature, efficiency, and effectiveness of the learning process. In this section, we systematically review the main methodologies for reward design, categorized by the nature of the feedback: textual feedback, internal confidence, external rewards, and implicit rewards.

\begin{figure}[tb]
  \centering
  \includegraphics[width=\linewidth]{Figures/how_reward.pdf}
  \caption{Overview of reward-based self-evolution strategies, categorized into textual, implicit, internal, and external rewards, each associated with distinct feedback sources and mechanisms.}
  \label{fig:how_reward}
\end{figure}




\paragraph{Textual Feedback}

Textual Feedback leverages the native modality of LLMs—natural language—to provide detailed, interpretable instructions for refinement. Unlike scalar rewards, textual feedback encapsulates nuanced critiques and actionable suggestions. Recent frameworks such as Reflexion \citep{shinn2023reflexion}, AdaPlanner \citep{sun2023adaplanner}, AgentS2 \citep{agashe2025agents2}, SELF \citep{lu2023self}, Self-Refine \citep{madaan2023self_refine}, SCoRe \citep{kumar2024training}, PAG \citep{jiang2025pag}, and TextGrad \citep{yellamraju2024textgrad} exemplify this direction. For instance, Reflexion proposes “verbal reinforcement learning,” where agents reflect in natural language on their past trials, storing these reflections as episodic memory to guide future decisions. AdaPlanner enables closed-loop adaptive planning by allowing LLM agents to revise their plans based on both in-plan and out-of-plan feedback, while also mitigating hallucination via code-style prompts and leveraging skill discovery. Self-Refine and SELF further explore iterative self-feedback and self-correction, demonstrating that even state-of-the-art models can be improved via multi-turn, language-based self-critique, without additional supervised data or external reinforcement. Such frameworks highlight the power of language as a reward channel, enabling nuanced, flexible, and sample-efficient self-improvement.


\paragraph{Internal Rewards}

Internal Confidence-based rewards move away from external signals and instead exploit internal metrics such as the model’s probability estimates or certainty. This paradigm leverages the model's intrinsic understanding to guide improvement without relying on external supervision. Methods such as Confidence-Informed Self-Consistency (CISC) \citep{taubenfeld2025cisc}, Self-Ensemble \citep{xu2025selfensemblemitigatingconfidencedistortion}, Self-Rewarding Self-Improving \citep{simonds2025selfrewardingselfimproving}, scalable best-of-N selection via self-certainty \citep{kang2025scalablebestofnselectionlarge}, and Self-Rewarding Language Models \citep{yuan2025selfrewardinglanguagemodels} allow models to self-evaluate and calibrate their responses based on internal confidence metrics. For example, CISC weights reasoning paths by confidence scores to improve both accuracy and computational efficiency, effectively filtering high-quality solutions from multiple candidates. Self-Ensemble mitigates confidence distortion by dividing choices into smaller, more manageable groups and aggregating predictions to reduce overconfidence bias. Self-Rewarding Language Models demonstrate that models can act as their own reward function, generating training data through self-instruction and self-evaluation cycles. These approaches can reduce reliance on human labels and external evaluators, enabling scalable and autonomous self-improvement loops that can operate continuously without human intervention. AgentEvolver~\citep{zhai2025agentevolver} proposes a comprehensive framework to improve agent training efficiency through three synergistic mechanisms: self-questioning for autonomous task generation, self-navigating for experience-guided exploration, and self-attributing for fine-grained credit assignment. In particular, its self-attributing mechanism uses an LLM's reasoning to retrospectively assign step-wise rewards that are dense and semantically grounded for policy optimization.

\paragraph{External Rewards}

External Rewards are derived from sources outside the model, such as the environment, majority voting, or explicit rules. Majority voting \citep{shafayat2025largereasoningmodelsselftrain, wei2025unsupervisedposttrainingmultimodalllm, zhang2025consistentpathsleadtruth} uses consensus among multiple model outputs as a proxy for correctness, providing a self-generated but grounded reward signal. Environment feedback, including tool-based signals, is central to agentic LLM research (e.g., SWE-Dev \citep{du2025swedevevaluatingtrainingautonomous}, SICA \citep{robeyns2025self}, Feedback Friction \citep{jiang2025feedbackfriction}, USEagent \citep{applis2025useagent}, DYSTIL \citep{wang2025dystildynamicstrategyinduction}), where agents learn through direct interaction with real-world environments and tools. Rule-based rewards \citep{wang2025otcpo, wang2025autorule, zhang2025learning, simonds2025ladder, wang2025ragen, liu2025spiral} use explicit constraints or logical rules as verifiable signals, particularly effective in the domains of mathematical reasoning, game play, and structured problem solving. These methods offer objective, reliable supervision but may require significant engineering or be limited in expressiveness.

\paragraph{Implicit Rewards}


Implicit Reward frameworks hypothesize that LLMs can learn from feedback signals even when not explicitly labeled as rewards. For instance, “Reward Is Enough” \citep{song2025reward} demonstrates that LLMs can perform in-context reinforcement learning using simple scalar signals embedded in the context window, improving their responses over rounds without explicit RL fine-tuning or supervision. This reveals an inherent capacity for models to interpret and learn from implicit feedback cues present in their input context. Recent work has expanded this concept by showing that LLMs inherently encode reward-like signals through their standard training objectives. 
Endogenous reward~\citep{li2025generalist} reveal that standard next-token prediction implicitly learns a generalist reward function, which can be extracted from model logits without additional training. Moreover, ImPlicit Self-ImprovemenT (PIT) framework~\citep{wang2024enabling} implicitly learns the improvement goal
from human preference data without extra human efforts by maximizing the quality
gap of the response conditioned on a reference response. Unlike rule-based or environment-derived external rewards, implicit reward methods offer unique advantages by discovering and utilizing reward signals that are inherently present in language modeling.


In summary, reward-based evolution provides explicit optimization and strong autonomy but remains sensitive to reward design, often trading stability and safety for adaptability and openness.



\subsection{Imitation and Demonstration Learning}


Imitation and demonstration learning traditionally involves an agent that learns to mimic the behavior of an expert (typically a human) from a set of demonstrations. In the context of self-evolving agents, this paradigm is adapted and generalized, which is the focus of our survey. Here, the role of the "expert" is not necessarily a fixed, external entity (e.g., human) but rather any source of high-quality demonstration. In self-evolving agents, these "expert exemplars" are typically generated by the agent itself (e.g., a past successful trajectory), by other more capable agents, or synthesized from environmental interactions.

The key distinction between imitation learning and reward-based methods lies in the nature of the feedback. Imitation learning is \textit{prescriptive} and \textit{exemplar-based}: the agent is provided with a complete, successful guide (e.g., a full reasoning trace) and learns to reproduce this behavior. In contrast, reward-based methods are \textit{evaluative} and \textit{signal-based}: the agent explores on its own and receives a scalar or textual critique, forcing it to infer the path to improvement through trial-and-error and credit assignment.

Furthermore, the evolutionary mechanism differs fundamentally from population-based methods that will be introduced later. While both imitation and reward-based learning typically focus on optimizing a \textit{single agent's improvement} through iterative refinement, population-based methods evolve a \textit{collection of agents} in parallel. Their progress typically comes from selection pressure across a diverse gene pool, rather than the direct knowledge transfer from an exemplar to an individual. Therefore, imitation learning occupies a unique niche: it relies on the availability of high-quality solutions to directly guide and accelerate the evolution of an individual agent, making it exceptionally powerful when such demonstrations can be reliably and autonomously generated.




\subsubsection{Self-Generated Demonstration Learning}

Self-generated demonstration learning involves agents creating their own training data through iterative refinement processes, where the models learn to improve by generating and selecting high-quality examples from their own outputs.

\textbf{Bootstrapping Reasoning Capabilities.} 
\cite{zelikman2022star} introduces the foundational framework for self-generated demonstration learning, enabling language models to bootstrap their reasoning capabilities through iterative self-training. This process involves generating reasoning chains for problems, fine-tuning on correct solutions, and repeating this cycle to progressively improve performance without the need for ground-truth reasoning paths. Building on this framework, recent advancements have refined the bootstrapping process through more sophisticated training strategies. For instance, \cite{hosseini2024v} proposes a verifier-guided self-training approach, where separate verifier models assess the quality of generated reasoning chains before they are incorporated into the training data, enhancing the reliability of self-improvement. Additionally, \cite{koh2025adastar} introduces adaptive data sampling strategies that dynamically adjust the composition of training data based on model performance across various reasoning tasks, thereby mitigating overfitting to specific problem types. The "Explore to Evolve" paradigm~\citep{wang2025exploretoevolve} extends this concept to deep research web agents by proposing an automated pipeline for generating complex, verifiable training data. The framework directs an agent to first perform proactive online exploration to gather grounded information from the live web, and then to evolve a sophisticated aggregation logic to synthesize question-answer pairs that require both information-seeking and deep reasoning.

\textbf{Multimodal Self-Training.} 
Extending self-training to multimodal domains presents unique challenges in generating high-quality demonstrations that span both visual and textual modalities. \cite{deng2024enhancing} demonstrates how vision-language models can improve iteratively by training on their own generated image descriptions and visual reasoning chains. The approach leverages the model’s existing visual understanding to generate detailed image descriptions, which are subsequently used to fine-tune the model’s visual perception in a bootstrapping manner. \cite{zhao2024genixer} builds on this concept by empowering multimodal large language models to serve as powerful data generators, producing diverse training examples across different modalities and tasks through advanced prompt engineering and quality filtering mechanisms.

\subsubsection{Cross-Agent Demonstration Learning}

Cross-agent demonstration learning involves agents learning from demonstrations provided by other agents, either within the same system or from external sources, enabling knowledge transfer and collaborative improvement.


\textbf{Multi-Agent Bootstrapped Reasoning.} 
\cite{zhao2025sirius} presents a framework for multi-agent systems to learn from each other’s successful demonstrations through bootstrapped reasoning. The system maintains an experience library containing successful interaction trajectories generated by different agents, facilitating efficient knowledge sharing and collaborative improvement. Each agent can leverage the collective experience of the entire system, thereby accelerating the learning process and enabling the discovery of diverse solution strategies. This framework illustrates how agents can specialize in different aspects of complex tasks while benefiting from the accumulated knowledge of the entire system.

\textbf{Domain-Specific Demonstration Learning.} 
Domain-specific applications of demonstration learning have proven especially effective in specialized fields where expert knowledge can be effectively transferred through demonstrations. In recommendation systems, techniques such as self-optimized fine-tuning \citep{tang2025bridging} enable LLM-based recommender systems to learn from their own successful recommendation patterns, creating a feedback loop that enhances personalization over time. The system generates high-quality recommendation demonstrations from successful user interactions and uses these to fine-tune the underlying language model, ultimately leading to more accurate and personalized recommendations.

\subsubsection{Hybrid Demonstration Learning}

Hybrid demonstration learning combines both self-generated and external demonstrations to create more robust and diverse training regimens that leverage the strengths of each approach.

\textbf{Recursive Self-Improvement.} 
\cite{qu2024recursive} demonstrates how agents can be trained to systematically improve their behavior through structured self-reflection and demonstration generation. This approach enables language model agents to introspect on their reasoning processes, identify areas for improvement, and generate corrective demonstrations to address these weaknesses. This recursive process establishes a continuous improvement loop, where agents become increasingly skilled at self-diagnosis and self-correction, leading to more robust and adaptable behavior.

\textbf{Confidence-Guided Demonstration Selection.} 
Recent developments have focused on more sophisticated mechanisms for selecting high-quality demonstrations from both self-generated and external sources. Confidence-based approaches \citep{li2024confidence} utilize the model's uncertainty estimates to determine which demonstrations are most likely to contribute positively to learning, filtering out potentially detrimental or low-quality examples. This method addresses a critical challenge in demonstration learning: poor-quality demonstrations can degrade performance. By ensuring that only high-confidence, high-quality examples are used for training, this approach helps to maintain the integrity of the learning process.

The effectiveness of imitation and demonstration learning approaches is highly dependent on the quality and diversity of the available demonstrations. While these methods can yield impressive results when high-quality exemplars are present, they face challenges in domains where good demonstrations are scarce or where the optimal behavior is not well-represented in the available data. Future research directions include developing more sophisticated demonstration selection and generation strategies, improving the robustness of learning from imperfect demonstrations, and creating better mechanisms for combining demonstrations from multiple sources.



Overall, imitation-based evolution stabilizes learning through high-quality exemplars but often trades exploration and generalization for reliability and sample efficiency.


\subsection{Population-based and Evolutionary Methods}


Population-based and evolutionary methods are a paradigm with a long history in improving agent behavior that complements modern learning-based techniques. This approach, drawing inspiration from biological evolution, has deep roots in AI. The concept was formalized into a practical computational tool by John Holland, whose seminal work on the Genetic Algorithm (GA) established the core operators of selection, crossover, and mutation for refining a population of solutions~\citep{holland1976adaptation}. Building on this, John Koza pioneered Genetic Programming (GP), a powerful extension that directly evolves executable programs or symbolic expressions, which makes it appropriate for generating agent logic. This paradigm's power was demonstrated by automatically synthesizing novel, human-competitive results, such as patented analog electrical circuits~\citep{koza2010human}. This success extended into diverse domains, from evolving competitive agents for strategic games like backgammon and chess~\citep{hauptman2007evolution, azaria2005gpgammon} to discovering complex, interpretable policies in agent-based simulations, such as evolving dynamic taxation rules that outperformed static, human-designed strategies~\citep{garuccio2016genetic}.

This paradigm led to landmark achievements in evolving agent systems. For example, the classic "Evolved Synthetic Creatures" co-evolved agent morphology and neural controllers in a simulated 3D world, leading to the discovery of a wide variety of novel and effective locomotion strategies~\citep{karl1994evolving}. Later, the influential NEAT algorithm addressed a critical challenge by demonstrating how to evolve not just the weights but the entire topology of a neural network~\citep{stanley2002evolving}. This enabled the autonomous discovery of complex agent "brains" from simple initial structures, a principle that has had an enduring impact on neuroevolution. These seminal works illustrated that evolution could construct both an agent's physical form and its complex control systems, establishing a powerful alternative to manual design.

Building on this foundation, these methods represent a different paradigm for agent evolution compared to the reward-based and imitation-based approaches discussed in previous sections. While reward-based methods typically optimize individual agents through iterative reward signals and imitation learning relies on learning from demonstrations, population-based methods maintain multiple agent variants simultaneously. This allows for parallel exploration of the solution space and the emergence of diverse capabilities through mechanisms such as selection, mutation, crossover, and competitive interaction \citep{zhang2025nature}. By leveraging parallel search and genetic variation, these methods enable broader search coverage and the discovery of novel solutions that might be missed by gradient-based optimization. This approach is particularly valuable when the solution space is complex, multimodal, or when the optimal strategy requires fundamental architectural changes rather than parameter fine-tuning.




\subsubsection{Single Agent Evolution}

Single-agent evolutionary approaches focus on evolving individual agents through population-based mechanisms, where multiple variants of an agent compete and evolve over time. These methods can be broadly categorized into two main paradigms: learning from evolution and self-play from multiple rollouts.

\textbf{Learning from Evolution.} 
This paradigm draws directly from biological evolution, maintaining populations of agent variants and applying evolutionary operators to discover improved capabilities. The Darwin Gödel Machine (DGM) \citep{zhang2025darwin} exemplifies this approach through open-ended evolution of self-improving agents that maintain an archive of all historical versions, enabling branching from any past "species" rather than linear optimization. The system achieves self-referential improvement by allowing agents to directly modify their own Python codebase, with evolution driven by empirical performance on coding benchmarks and parent selection balancing performance scores with novelty rewards for diverse exploration. Recent work has further explored using LLMs themselves to implement core evolutionary operators. For instance, LLM\_GP~\citep{hemberg2024evolving} uses the LLM to perform mutation, crossover, and selection directly on programs represented as text, leveraging the model's innate knowledge of code to inform the evolutionary search. Similarly, open-source frameworks like CodeEvolve~\citep{assumpccao2025codeevolve} have demonstrated that evolutionary coding agents can achieve state-of-the-art results on mathematical benchmarks, sometimes outperforming proprietary systems like AlphaEvolve~\citep{novikov2025alphaevolve} by using an island-based genetic algorithm and inspiration-based crossover. This principle is also explored in Self-Referential Graph HyperNetworks~\citep{pedersen2025hypernetworks}, where networks learn to generate their own weight mutations, allowing the rate of evolution itself to become selectable and adaptable.

Beyond evolving code and architecture, this paradigm extends to evolving the model's parameters and internal logic. The Nature-Inspired Population-Based Evolution (GENOME) framework \citep{zhang2025nature} directly applies genetic algorithms to language model parameter evolution, maintaining populations and using crossover, mutation, and selection operators on model weights. GENOME+~\citep{zhang2025nature} extends this with particle swarm optimization concepts, adding inheritance mechanisms and ensemble methods that demonstrate gradient-free evolutionary optimization can effectively improve model capabilities through parameter space exploration. EvoLLM-JP~\citep{akiba2025evolutionary} takes this further by using evolutionary algorithms to optimally merge multiple foundation models into a single, specialized model with superior performance. Furthermore, some frameworks create a tightly integrated feedback loop where evolution helps model fine-tuning. SOAR~\citep{pourcel2025soar}, for example, alternates between an evolutionary search phase to generate candidate programs and a "hindsight learning" phase that uses all attempts (both successful and failed) to generate a rich dataset for fine-tuning the agent model, creating a virtuous cycle of self-improvement.




\textbf{Self-Play.}
Self-play is a paradigm where agents improve through iterative interaction with versions of themselves, creating a dynamic and self-sustaining learning process. Its principles were famously demonstrated by systems like AlphaZero~\citep{silver2017mastering}, which achieved superior performance in complex games by learning entirely without human data. The core mechanism is co-evolutionary learning: as an agent improves, its opponents (past or concurrent versions of itself) also become stronger, generating a perpetual and adaptive curriculum of increasing difficulty. This avoids the stagnation that can occur when training against a fixed environment and enables the discovery of novel, emergent strategies.

This powerful principle has been adapted for LLMs and LLM Agents, enabling them to bootstrap their capabilities from zero or minimal external data. A prominent approach involves a single model or two model instances adopting distinct, co-evolving roles. For instance, Absolute Zero \citep{zhao2025absolute} and R-Zero \citep{huang2025r} employ a "challenger" or "proposer" agent that generates problems at the frontier of a "solver" agent's capabilities. A more complex multi-agent dynamic is seen in Socratic-Zero \citep{wang2025socratic}, where a Solver co-evolves with a powerful Teacher that creates challenges and a Generator that distills the Teacher's strategy for scalable curriculum creation. To address the instability of purely autonomous systems, R-Few \citep{yu2025guidedself} introduces a guided approach where the challenger is grounded by a small set of human examples to prevent concept drift and diversity collapse.
The system improves through a closed loop where the solver is rewarded for correctness (often verified by execution) and the challenger is rewarded for posing difficult yet solvable problems, thus driving continuous improvement without external labels. 
Similarly, Self-Challenging Language Model Agents \citep{zhou2025selfchallenginglanguagemodelagents} establishes a framework where an agent alternates between generating and solving complex, multi-step coding tasks, using successful trajectories to fine-tune itself. The paradigm also extends to more specialized, collaborative roles, as seen in the Sol-Ver framework \citep{lin2025learningtosolve}, where an LLM co-evolves its ability to both generate code (solver) and create corresponding unit tests (verifier). Likewise, SPELL~\citep{yang2025spell} applies this principle to long-context reasoning, with a single model cyclically adopting questioner, responder, and verifier roles to provide reliable reward signals in a domain where programmatic verification is difficult.

Across these approaches, improvement is driven by self-generated learning signals derived from the agent’s own trajectories. Self-Play Fine-Tuning (SPIN) \citep{chen2024self} establishes a foundational approach where current models compete against previous versions, creating evolutionary pressure for improvement. SPC \citep{chen2025spc} advances this with a more sophisticated adversarial co-evolution, featuring a "sneaky generator" that creates deceptive errors and a "step critic" that learns to detect them. STL \citep{mendes2025language} demonstrates self-teaching through iterative lookahead search, where value models generate training data from their own exploratory rollouts. Recent work, such as EvoTest \citep{he2025evotest}, extends these ideas by introducing a gradient-free, evolutionary framework that revises an agent’s prompt, memory, and tools between episodes.

Self-play is distinguished by a unique self-improvement mechanism: an agent learns through direct interaction with variations of itself. This typically manifests in two ways: a model competes against its own past versions to drive iterative refinement (as in SPIN), or a single model adopts distinct, interacting roles, such as a "challenger" generating novel problems for a "solver" (as in Absolute Zero). This principle of learning from \textit{dynamic interaction} is different from imitation-based bootstrapping. While methods like STaR~\citep{zelikman2022star} also learn from an agent's own outputs, they do so by filtering and training on static, successful trajectories. Self-play, in contrast, generates its learning signal from the process of the game-like interaction itself, learning from relative success even when no perfect exemplar exists. This focus on a single agent's lineage also sets it apart from broader population-based methods: instead of evolving a large, diverse population, self-play creates a highly focused evolutionary pressure between a minimal set of policies derived from the same agent.



\subsubsection{Multi-Agent Evolution}

Multi-agent evolutionary methods extend population-based approaches to evolving entire teams or networks of agents, focusing on optimizing collective behavior, coordination strategies, and collaborative architectures. These approaches can be categorized into two main paradigms based on their evolution mechanisms: System Architecture Evolution and Knowledge-Based Evolution.

\textbf{System Architecture Evolution.}
This paradigm focuses on evolving the structural and coordination aspects of multi-agent systems, including team composition, orchestration strategies, and workflow optimization. 
EvoMAC \citep{hu2024self} introduces a framework that mimics neural network training for multi-agent systems, implementing "textual backpropagation" where compilation errors and test failures serve as loss signals to drive iterative modifications of agent team composition and individual prompts. A specialized "updating team" analyzes textual feedback to identify problematic agents and generate modification instructions, effectively implementing gradient-based optimization in the space of agent configurations rather than model parameters. The FELA framework~\citep{wang2025fela} applies this concept to a practical industrial problem, using a multi-agent system with specialized "Idea," "Code," and "Critic" agents that collaboratively evolve to generate high-performing features from complex data, guided by principles from both reinforcement learning and genetic algorithms. Puppeteer \citep{dang2025multi} takes a different approach by focusing on coordination strategy evolution rather than team composition changes. The system employs a centralized orchestrator that evolves its decision policy through reinforcement learning, dynamically selecting which agents to activate at each step while balancing task performance with computational cost. This "puppeteer-puppet" paradigm demonstrates how architectural evolution can occur at the coordination level, discovering efficient collaboration patterns and emergent behaviors such as tighter coordination among core agents and sophisticated cyclic interaction patterns. Agent0~\citep{xia2025agent0} introduces a framework that evolves agents from zero data via a co-evolutionary loop between a curriculum agent and an executor agent. The curriculum agent is trained to generate frontier tasks that challenge the executor. Then, the improved tool-use capabilities of the executor in turn drive the creation of a more complex, tool-aware curriculum, establishing a virtuous cycle of self-improvement.

\textbf{Knowledge-Based Evolution.}
This paradigm emphasizes evolving the collective knowledge and experience of multi-agent teams through memory accumulation and case-based learning, primarily operating through in-context learning or in-context-like adaptation rather than parameter updates. MDTeamGPT \citep{chen2025self} establishes the foundation for this approach through a dual knowledge base system, implementing CorrectKB for storing successful cases and ChainKB for capturing failure reflections, enabling the system to learn from both successes and mistakes through structured case retrieval and reasoning enhancement. Extending this medical consultation framework, MedAgentSim \citep{almansoori2025self} demonstrates how such knowledge-based evolution can be applied to real-world diagnostic scenarios, accumulating experience from patient interactions and using retrieval-augmented generation to improve consultation quality over time. 
PiFlow \citep{pu2025piflow} applies this paradigm to scientific discovery, maintaining a trajectory of principle-outcome pairs and using them to steer hypothesis generation through information-theoretical optimization.


In summary, population-based and self-play evolution enhance diversity and open-ended discovery, yet typically incur higher computational cost and lower interpretability compared with single-agent paradigms.





\subsection{Cross-cutting Evolutionary Dimensions}


After outlining the three core evolutionary paradigms, we now analyze their cross-cutting dimensions—revealing how different design choices balance feedback type, data source, and learning stability.


\begin{figure}[t]
  \centering
  \includegraphics[width=0.9\linewidth]{Figures/how_cross_dim.pdf}
  \caption{
Illustration of the cross-cutting evolutionary dimensions underlying agent self-evolution. 
\textbf{Learning Paradigm:} in \textit{offline learning}, data are pre-collected and used for filtering and fine-tuning, while in \textit{online learning}, agents continuously interact with their environments for real-time adaptation. 
\textbf{Policy Consistency:} \textit{on-policy evolution} updates policies based on the agent’s own trajectories, whereas \textit{off-policy evolution} relies on replay buffers, human demonstrations, or experiences from other agents. 
\textbf{Reward Granularity:} feedback can be \textit{process-based} (step-level rewards), \textit{outcome-based} (final-result rewards), or a \textit{hybrid} of both. 
Together, these three orthogonal axes define how agents generate data, adapt their policies, and receive feedback across reward-based, imitation-based, and population-based evolution paradigms.
}

  \label{fig:cross}
\end{figure}


Agent self-evolution is a multifaceted process characterized by a number of cross-cutting dimensions that shape how agents learn, adapt, and improve over time. Beyond any single learning algorithm or supervision signal, these dimensions define the core principles underlying the design and analysis of autonomous agents. In this section, we systematically compare the major families of self-evolution methods—reward-based, imitation/demonstration-based, and population-based—along several key axes, such as learning paradigm (online vs. offline), policy consistency (on-policy vs. off-policy), and reward granularity (process-based, outcome-based, or hybrid). We further highlight additional dimensions, including feedback types, data sources, sample efficiency, stability, and scalability, as summarized in Table~\ref{tab:comparison_dimension}. This comprehensive comparison provides a unified perspective for understanding the strengths, limitations, and design trade-offs inherent in different approaches to agent evolution.


\begin{table}[t]
\centering
\caption{Comparison of self-evolution method families along key dimensions.}
\label{tab:comparison_dimension} %
\rowcolors{2}{gray!10}{white} %
\begin{tabularx}{\textwidth}{l XXX} %
\toprule
\textbf{Dimension} & \textbf{Reward-based} & \textbf{Imitation/Demonstration} & \textbf{Population-based} \\
\midrule
\textbf{Feedback Type} 
  & Scalar reward, natural language, confidence, external signals
  & Demonstration trajectories, exemplars, rationales
  & Fitness scores, task success, competitive signals \\
\textbf{Data Source} 
  & Self-generated, environment, external rules
  & Self-generated or other agents, humans
  & Population generations, multi-agent systems \\
\textbf{Reward Granularity} 
  & Outcome/process/hybrid (flexible)
  & Usually outcome/process (via demo steps)
  & Often outcome-level, sometimes process via competition \\
\textbf{Online/Offline} 
  & Both (reward learning, RL, DPO, SFT)
  & Typically offline, sometimes online demo mining
  & Online evolution or batch population updates \\
\textbf{On/Off-policy} 
  & Both (DPO, Reflexion, GRPO)
  & Primarily off-policy, but online variants can be on-policy
  & Off-policy (population); self-play is on-policy\\
\textbf{Sample Efficiency} 
  & Moderate (depends on reward sparsity)
  & High (if demo quality is high)
  & Usually low (needs many trials) \\
\textbf{Stability} 
  & Sensitive to reward design
  & Sensitive to demo quality/diversity
  & Sensitive to population size/diversity\\
\textbf{Scalability}
  & Good with automation
  & Limited by demo collection
  & High but resource-intensive \\
\bottomrule
\end{tabularx}
\end{table}



\subsubsection{Online and Offline Learning}
Another fundamental dimension in the design of self-evolving agents is the learning paradigm, which can be broadly categorized as either offline or online. This distinction depends on whether the agent’s evolutionary updates are performed on a static, pre-collected dataset of experiences (offline) or through continuous, direct interaction with a live environment (online).

\textbf{Offline Learning}
In the offline learning paradigm, the learning phase is decoupled from live task execution. The offline process typically involves cycles of offline data generation, filtering, and model fine-tuning, focusing on building a powerful and generalist foundational model before deployment.
A primary strategy in this domain is LLM bootstrapping, where a model enhances its own capabilities using its self-generated content. For example, Self-Instruct\citep{wang2022self} shows how a language model can bootstrap its own instruction-following ability by generating new instructions, paired with its own responses, creating a synthetic dataset for fine-tuning. Building on this, WizardLM\citep{xu2024wizardlm} demonstrates how to progressively evolve the complexity of these self-generated instructions, pushing the model’s capabilities on more challenging tasks. Although these methods primarily focus on broad capability expansion via synthetic heuristics, acting as a bootstrapping phase, they lay the necessary groundwork for closed-loop, experience-driven evolution defined in our framework.
In the context of GUI and Web agents, offline learning often involves leveraging pre-collected high-quality trajectories for supervised fine-tuning (SFT). OS-Genesis\citep{sun2024genesis} introduced a reverse task synthesis method for automatic trajectory creation. Similarly, UI-Genie\citep{xiao2025ui} employs a unified reward model for trajectory evaluation and a self-improving loop to generate high-quality trajectories iteratively. Both approaches focus on curating a rich SFT dataset to enhance the agent’s capabilities to solve complex tasks.
Beyond SFT, offline methods also incorporate reinforcement learning performed on a static dataset of agent-environment interactions. For example, GUI-R1\citep{luo2025gui} and InfiGUI-R1\citep{liu2025infigui} utilize rule-based rewards and apply R1-style\citep{guo2025deepseek} training on offline GUI datasets.

\textbf{Online Learning}
In contrast, online learning enables an agent to learn and adapt continuously while it interacts with a live or simulated environment. Feedback from each action is used to update the agent’s policy, plan, or knowledge base in real-time. This allows for greater adaptability to dynamic or unseen situations.
Some agents evolve online not by updating their model weights, but by refining their plans and skill libraries on the fly. For example, Voyager\citep{wang2023voyager} presents an LLM-powered agent that learns to play Minecraft by continuously exploring, generating its own curriculum of tasks, and building a persistent skill library from direct experience.
AdaPlanner\citep{sun2023adaplanner} focuses on adapting its plan within a task; it generates an initial plan, receives feedback from the environment, and refines the plan online. Similarly, SwiftSage\citep{lin2023swiftsage} operates with a fast-and-slow thinking process, where it can reflect on failures of its fast, intuitive mode and switch to a more deliberate, tool-using slow mode, adapting its strategy online based on task difficulty.
Reinforcement Learning serves as a fundamental mechanism for online learning, enabling agents to learn from environmental reward signals. DigiRL\citep{bai2024digirl} demonstrates how to train device-control agents in the wild using autonomous RL, while DistRL\citep{wang2024distrl} proposes an asynchronous distributed framework to make such on-device training feasible. MobileGUI-RL\citep{shi2025mobileguirladvancingmobilegui} addresses the specific challenges of training GUI agents in online mobile environments by introducing a synthetic task generation pipeline combined with group relative policy optimization (GRPO) through trajectory-aware rewards.

\subsubsection{On-policy and Off-policy Learning}

While the previous section examined the timing of data collection and learning (online vs offline), this section focuses on the policy consistency aspect of agent evolution - specifically, whether agents learn from experiences generated by the same policy they are trying to improve (on-policy) or from experiences generated by different policies (off-policy). This distinction is crucial for understanding how agents utilize their experiential data and manage the trade-offs between learning stability and sample efficiency during the evolutionary process.

\textbf{On-policy Learning.}
On-policy approaches require agents to learn exclusively from experiences generated by their current policy, ensuring policy consistency but often at the cost of sample efficiency. Reflexion \citep{shinn2023reflexion} exemplifies this approach through its iterative self-reflection mechanism. The agent generates responses using its current policy, receives feedback on failures, and immediately incorporates this feedback to update its reasoning process for the next iteration. GRPO \citep{shao2024deepseekmath} and DAPO \citep{yu2025dapo} continue this path and show the effectiveness of multiple rollouts. The agent always learns from its current behavior, maintaining strict policy consistency. In agent settings, on-policy methods provide excellent learning stability and avoid distribution mismatch issues that plague off-policy methods. However, they suffer from low sample efficiency, as each policy update requires fresh data collection, making them computationally expensive for complex multi-step reasoning or tool use scenarios where generating high-quality trajectories is costly.

\textbf{Off-policy Learning.}
Off-policy approaches allow agents to learn from experiences generated by different policies, including previous versions, other agents, or human demonstrations, significantly improving sample efficiency at the cost of potential distribution mismatch. 
 \cite{yuan2401self} demonstrates a sophisticated off-policy approach where model $M_{t+1}$ learns from preference data generated by the previous version $M_t$. The system handles distribution shift through DPO's built-in KL divergence constraint with the reference policy, preventing the new policy from deviating too far from the data-generating policy. 
\cite{yuan2023rrhf} showcases another powerful off-policy paradigm by learning from diverse response sources—including other models, humans, and different sampling strategies—through ranking-based supervision. The method elegantly sidesteps distribution shift by treating alignment as a ranking problem rather than requiring policy consistency. 
\cite{zhao2025sirius} illustrates off-policy learning in multi-agent settings, where agents learn from an "experience library" containing successful interaction trajectories generated by previous policy versions, enabling efficient reuse of expensive multi-agent coordination data. 
In agent settings, off-policy methods excel in sample efficiency, allowing agents to leverage historical data, expert demonstrations, and cross-agent learning. They are particularly valuable for multi-step reasoning where successful trajectories are rare and expensive to generate, and for tool use scenarios where agents can learn from diverse execution examples without repeated environmental interaction. However, they face challenges with distribution shift, reward hacking (where agents exploit inconsistencies between training and deployment policies), and the need for careful regularization to maintain training stability.

\subsubsection{Reward Granularity}

Another critical choice in the reward design is its granularity, which determines at what level of detail the agent receives its learning signal.
Reward granularity ranges from coarse-grained outcome-based rewards, which evaluate the overall task completion, to fine-grained process-based rewards that assess each step of the agent’s trajectory. Current self-evolution frameworks adopt these varying levels of granularity to tailor feedback mechanisms according to task complexity and the desired learning outcomes.

\textbf{Outcome-based Reward}
Outcome-based Reward is a feedback mechanism that evaluates an agent based on the successful completion of predefined tasks. This reward is determined solely by the final state of the agent’s trajectory, regardless of the intermediate steps. 
A central challenge, particularly in dynamic environments like web or GUI navigation, is to effectively learn from both successful trajectories and the much more frequent failure trajectories.
To address this, Direct Preference Optimization (DPO)\citep{rafailov2023direct} is designed to directly maximize the likelihood of preferred responses while minimizing the KL-divergence with the reference policy.
Similarly, RRHF\citep{yuan2023rrhf} employs a ranking loss approach that aligns model probabilities of multiple responses with human preferences by ranking response probabilities without requiring auxiliary value models.
Moreover, several works have developed specialized frameworks for agent self-evolution that are built upon outcome-based rewards.
A straightforward approach is rejection sampling finetuning, as used in AutoWebGLM\citep{lai2024autowebglm}. This method employs a pre-designed reward model to evaluate trajectory outcomes, identify the successful trajectories, and update the model with this high-quality data.
DigiRL\citep{bai2024digirl} models the GUI navigation task as a Markov Decision Process (MDP) and obtains a final, sparse reward at the end of an episode using a VLM-based evaluator.
WebRL\citep{qi2024webrl} develops a robust outcome-supervised reward model (ORM) to address the feedback sparsity inherent in dynamic web environments. The ORM evaluates task success within a self-evolving curriculum framework, enabling agents to learn from unsuccessful attempts and progressively improve.

      
\textbf{Process-based Reward}
In contrast to outcome-based rewards, which provide a single, delayed signal, the process-based reward paradigm offers more precise and granular feedback by evaluating each step in an agent’s trajectory.
Process-supervised reward models (PRMs) have been demonstrated to be significantly more reliable than outcome-supervised reward models (ORMs), particularly in domains requiring complex reasoning like solving math problems\citep{lightman2023let}.
However, obtaining such fine-grained step-level feedback traditionally requires extensive human annotations, which are both time-consuming and expensive to scale.
To address this annotation bottleneck, Math-Shepherd\citep{wang2023math} proposes an automatic process annotation framework that utilizes Monte Carlo Tree Search (MCTS) to gather step-wise supervision by assessing each step’s potential to derive the correct final answer. Similarly, AlphaMath\citep{chen2024alphamath} trains a value model to evaluate the step correctness in solution paths and updates both the policy and value model through exploration and exploitation within an MCTS framework.
By leveraging process-based rewards, agents can improve their capabilities in a progressive, step-by-step manner.
rStar-Math\citep{guan2025rstar} and AgentPRM\citep{choudhury2025process} both propose methods to iteratively evolve the policy and the process reward model, generating progressively higher-quality reasoning paths without manual labels. Agent Q\citep{putta2024agentq} integrates a step-wise verification mechanism into its MCTS process to collect high-quality trajectories, which are then used to iteratively refine the policy via DPO training. 

\textbf{Hybrid Reward}
The hybrid methods aim to provide more comprehensive learning signals by incorporating both the clarity of final task success (outcome-based) and the granular guidance of intermediate steps (process-based). These methods overcome the sparsity of outcome-only signals while grounding the agent’s step-by-step reasoning in the ultimate task goal.
For example, GiGPO\citep{feng2025group} addresses the instability of training long-horizon agents by introducing a dual-level reward mechanism. It provides an episode-level reward based on the final success of entire trajectories, while simultaneously assigning a localized, step-level reward for intermediate actions. This dual signal provides both a high-level directional goal and low-level corrective guidance.
Similarly, SPA-RL\citep{wang2025spa} proposes a reward decomposition method that bridges the gap between sparse outcome signals and dense process feedback. It attributes incremental progress to each step within multi-step trajectories based on the final task completion, effectively distributing the outcome-based reward across the process steps. This approach creates dense intermediate progress rewards that enhance reinforcement learning effectiveness while maintaining alignment with the ultimate task objectives.







 

\subsection{Other Dimensions of Self-Evolution Methods}

In addition to the core axes of learning paradigm, policy consistency, and reward granularity, Table~\ref{tab:comparison_dimension} highlights several other important dimensions that differentiate self-evolution methods:

\textbf{Feedback Type.}
The nature of feedback varies widely: reward-based methods leverage scalar rewards, natural language signals, or model confidence; imitation methods focus on demonstration trajectories and rationales; population-based methods use fitness scores or competitive signals. The feedback type fundamentally determines what information the agent uses to improve.

\textbf{Data Source.}
Reward-based methods typically generate data through agent-environment interaction or engineered rules, while imitation learning often relies on human or expert-generated demonstrations. Population-based approaches draw from the collective experience of multiple agents or generations, enabling diverse exploration but requiring significant coordination.

\textbf{Sample Efficiency.}
Imitation learning is generally the most sample-efficient, provided high-quality demonstrations are available, as agents can directly mimic expert behavior. Reward-based methods are moderately efficient, with efficiency highly sensitive to reward sparsity. Population-based evolution tends to be sample-inefficient, as it often requires evaluating a large number of agent variants through many trials.

\textbf{Stability.}
Reward-based learning is sensitive to the quality and design of reward functions, risking reward hacking or unintended behaviors. Imitation learning depends heavily on the quality and diversity of demonstrations. Population-based methods are sensitive to population size and diversity, with small or homogeneous populations at risk of premature convergence.

\textbf{Scalability.}
Scalability is determined by the feasibility of data or feedback collection and the ability to parallelize learning. Reward-based methods scale well when feedback is automated (e.g., via simulators). Imitation learning is often bottlenecked by the cost of collecting demonstrations. Population-based approaches can scale to large compute but are highly resource-intensive.

Together, these dimensions offer a more nuanced, multidimensional view of self-evolution strategies, guiding practitioners in selecting and designing agent learning pipelines that are best matched to the challenges of their specific domains.
